\clearpage
\SourceLeaf{274}%
\begin{center}
\textbf{[METAPHYSICS OF GEOMETRY.]}

\medskip
[\textsc{W. Sartorius von Waltershausen.} \emph{Gauss zum Gedächtniss. Leipzig 1856. pp. 80---81.}]
\end{center}

\{\mbox{.\,.\,.\,.\,.\,.\,.\,.\,.} In his earliest youth, geometry had held little interest for him [\textsc{Gauss}], and that interest had developed to a high degree only later.

It was particularly remarkable and highly instructive to see the foundations on which mathematics rests laid bare by \textsc{Gauss} and sharply distinguished from metaphysics. Although he never published anything on these questions, there is reason to suppose that some material on them will be found among his scientific papers. Earlier, when the course of his life was still undecided and he had to reckon with the possibility of teaching mathematics somewhere, he had prepared a paper with that prospect in mind. It is said to have survived into his last years, and in it he had developed the beginnings of mathematics philosophically. Whether it can still be found is doubtful.

\textsc{Gauss} regarded geometry as a coherent structure only once the theory of parallels was admitted as its leading axiom. He had come to believe that this proposition could not be proved. Experience, however, showed it to be approximately correct, for example the angles of the triangle Brocken, Hohehagen, Inselsberg. If one declined to admit this axiom, the result was another, entirely\space%
